\documentclass[../../main.tex]{subfiles}
\begin{document}

\begin{center}
  {\fontsize{48pt}{58pt}\selectfont \textbf{1980}\normalsize\textbf{年度}}
\end{center}

\begin{table}[htb]
  \centering
  \caption{出題テーマと難易度}
  \begin{tabular}{|c|p{9cm}|c|}
    \hline
    \textbf{問題番号} & \textbf{テーマ} & \textbf{難易度} \\
    \hline
    第1問 & 二つの行列の積で得られる行列 & \\
    \hline
    第2問 & 放物線と座標軸に囲まれる面積の最大化 & \\
    \hline
    第3問 & 循環置換された3点と四面体 & \\
    \hline
    第4問 & 得点が0で止まるサイコロゲーム & \\
    \hline
    第5問 & 差について閉じた有限集合 & \\
    \hline
    第6問 & 整数係数の四次式の因数分解 & \\
    \hline
  \end{tabular}
\end{table}

\begin{table}[htb]
  \centering
  \caption{解答}
  \renewcommand{\arraystretch}{1.4}
  \begin{tabular}{|c|c|p{9cm}|}
    \hline
    \textbf{問題番号} & \textbf{小問} & \textbf{答え} \\
    \hline
    第1問 & -- &  \\
    \hline
    第2問 & -- & $b=4$のとき最大面積$\dfrac{25}{6}$ \\
    \hline
    第3問 & (1) & $A,B,C$は定円上にある \\
    \cline{2-3}
    & (2) & 四面体$OABC$の体積は$\dfrac{1}{6}$ \\
    \cline{2-3}
    & (3) & 指定された6点は定球面上にある \\
    \hline
    第4問 & -- & $5$回振れる確率は$\dfrac{5}{8}$，終了時の得点の期待値は$2$ \\
    \hline
    第5問 & (1) & 要素はすべて非負，またはすべて非正 \\
    \cline{2-3}
    & (2) & 適当に並べ替えると等差数列になる \\
    \hline
    第6問 & (i) & $g(0)=h(0)$ \\
    \cline{2-3}
    & (ii) & 両方が非定数なら$g=h$ \\
    \cline{2-3}
    & (iii) & 例：$(a,b)=(-1,-2)$ \\
    \hline
  \end{tabular}
\end{table}

\paragraph{試験形態}
6問．試験時間は要確認．

\end{document}
